\section{A.6. The anabelian sheets and the presentation calculations}
\subsection{The statements on 67/048--049 and the necessary choice of fundamental group}

We checked the original images, as well as the transcription. Sheet 049(a) asks to extend the anabelian conjecture to stable singular curves supplied with a normal propagator. Sheet 049(b) asks to recover nodes and marked points from conjugacy classes of procyclic subgroups invariant under the outer arithmetic action, and suggests that the action on these subgroups is automatically cyclotomic. Sheet 049(c) asks for a definition of an anabelian compactification in arbitrary dimension, together with a specialization theory. The marginal NB on 048 asks for three conjugacy classes of homomorphisms $\widehat{\Z}(1)\to\pi_1(X)$, with one class empty when there are no punctures. These are different questions; in particular, a theorem reconstructing a stable graph does not by itself reconstruct its component curves or its normal tensors.

We specify the logarithmic interpretation of the notation $[Y,N]$ before giving the reconstruction result. Let $K$ be a field finitely generated over $\Q$, let $Y$ be a geometrically connected stable marked curve of type $(g,n)$, $2g-2+n>0$, and let
\[
 N=(q_e)_{e\in\operatorname{Node}(Y_{\overline{K}})},
 \qquad 0\ne q_e\in T_{e^+}\widetilde{Y}\otimes
                         T_{e^-}\widetilde{Y}
\]
be Galois-equivariant. The two branches are unordered, and the tensor is interpreted using the symmetry of tensor product. On a split node, coordinates $x,y$ express the tensor as $q_e=c_e\,\partial_x\otimes\partial_y$. The local logarithmic smoothing is $xy=c_e t$ over the \emph{framed} standard log point: its chart generator $t$ is fixed. Equivalently, in the normal torus of the corresponding boundary stratum of the moduli stack, we take the $K$-rational tangential point with coordinates $(c_e)$.

The tangential construction can be defined by the pro-Kummer log chart: adjoining compatible roots of each normal parameter gives a pro-fs log scheme whose log fundamental group equals the ordinary fundamental group of the stratum point. The resulting map to the ambient tame fundamental group is the tangential map of \cite[Proposition 1]{HoshiMatsumotoNakayama2026A567}. Its Proposition 2 and Remark 1 show that changing a normal parameter by a unit congruent to $1$ along the stratum does not change this functor, up to its canonical homotopy. Thus first-order normal data suffice.

Let $\Delta_{Y,N}$ denote the geometric admissible, or logarithmic, fundamental group, equipped with the outer action
\[
 \rho_{Y,N}:G_K\longrightarrow\operatorname{Out}(\Delta_{Y,N})
\]
obtained from this tangential point. The underlying geometric group is the profinite surface group of type $(g,n)$. Its semi-graph of groups has as vertex groups the geometric fundamental groups of the normalization components with all marked points and node branches removed, and as edge groups the branch inertia groups, identified inversely at each node. This is the PSC construction of Mochizuki \cite[Definition 1.1 and Example 2.5]{Mochizuki2007A567}. We identify two maps of geometric groups only modulo inner automorphisms of the target.

The frame is material. Replacing $t$ by $a t$ changes the coefficients to $c_e/a$ and changes the arithmetic section by the Kummer class of $a$. It is therefore not an operation preserving the displayed $G_K$-action while forgetting its effect on the section. If one instead allows isomorphisms of the log base and discards its section, the data have been weakened. If one means ordinary $\pi_1^{\mathrm{et}}(Y)$, the normal tensor is absent altogether.

\begin{proposition}[Refuted for the ordinary fundamental group]\label{A6-ordinary-counterexample}
The ordinary arithmetic fundamental group of a stable marked curve cannot reconstruct a chosen nonzero normal tensor, even for $(g,n)=(0,4)$ over $\Q$. Moreover its geometric version does not retain nodal inertia.
\end{proposition}
\begin{proof}
Take two copies of $\mathbf P^1_\Q$, each with coordinates $z$, respectively $w$, glue their zeroes, and mark $1,\infty$ on each component, with four distinct labels. Every automorphism of this marked nodal curve is the identity, since each normalization component has three individually fixed points. Thus
\[
 (Y,\partial_z\otimes\partial_w),\qquad
 (Y,2\partial_z\otimes\partial_w)
\]
are not isomorphic, whereas their ordinary arithmetic fundamental groups are the same augmented group $\pi_1^{\mathrm{et}}(Y\setminus Y_\infty)\to G_\Q$: its definition makes no reference to either tensor. Its geometric group is the free profinite product of the ordinary groups of the two twice-punctured normalization components, hence $\widehat{F}_2$. To see this directly, a finite étale cover is a cover of each component, unramified also at the node, together with a bijection of their finite fibres at the node. There is no node ramification group in this description. In contrast, the admissible group glues the two triply punctured component groups along their nontrivial node inertia and is $\widehat{F}_3$, the group of a four-punctured smooth sphere. This proves both assertions and the necessity of the logarithmic distinction.
\end{proof}

\begin{lemma}[Proved: Kummer separation for the fields on 049]\label{A6-kummer}
For a field $K$ finitely generated over $\Q$, the Kummer map
\[
 K^\times\longrightarrow H^1(G_K,\widehat{\Z}(1))
\]
is injective. Also $(\widehat{\Z}(1))^{G_K}=0$.
\end{lemma}
\begin{proof}
Its kernel is $\bigcap_{m\ge1}(K^\times)^m$, by applying the usual Kummer exact sequence at every level $m$. Let $a$ lie in that intersection. Let $F$ be the algebraic closure of $\Q$ in $K$, a number field, and choose a normal projective model of $K/F$. Every divisorial valuation of $a$ is divisible by every positive integer and therefore vanishes. Normality implies that $a$ and $a^{-1}$ are global regular functions, so properness gives $a\in F^\times$. If $b^m=a$ with $b\in K$, then $b$ is algebraic over $F$, and hence $b\in F$. All finite-prime valuations of $a$ vanish, and the same holds for each such root $b$. Thus $a$ lies in every power subgroup of $\mathcal O_F^\times$. This group is finitely generated by the unit theorem \cite[Theorem 5.1]{MilneANTA567}; in a finitely generated abelian group the intersection of all power subgroups is the identity. Hence $a=1$.

For the second assertion, for each prime $\ell$ the cyclotomic character of $G_K$ has open image in $\Z_\ell^\times$: the intersection of $K$ with the cyclotomic algebraic extension of $\Q$ is contained in the number field $F$. Choose an element whose $\ell$-adic cyclotomic value differs from $1$. A fixed vector in $\Z_\ell(1)$ is killed by that nonzero difference, and so is zero. Applying this to every $\ell$ proves the assertion.
\end{proof}

\begin{lemma}[Known with derivation: graphical fullness over a finitely generated field]\label{A6-graphic}
For the framed logarithmic realization above, every $G_K$-compatible outer isomorphism
$\Delta_{Y,N}\simeq\Delta_{Y',N'}$ preserves the entire semi-graph of anabelioids: vertices, nodes, cusps, and incidence.
\end{lemma}
\begin{proof}
We verify the hypothesis of \cite[Corollary 2.7(ii)]{Mochizuki2007A567}, with $\Sigma$ the set of all primes and with any fixed $\ell$. The cyclotomic image is open by Lemma~\ref{A6-kummer}. For every characteristic finite admissible cover, the quotient denoted
$M^{\mathrm{vert}}/M^{\mathrm{edge}}$ in \cite[Definition 2.3(iii)]{Mochizuki2007A567} is the direct sum of the Tate modules of the Jacobians of its compact normalization components. The cover and its components descend to a finite extension of $K$. Spread these finitely many curves and Jacobians over a nonempty open of an integral scheme of finite type over $\Z[1/\ell]$; shrink until the component Jacobians are abelian schemes. A closed point of this open has a finite residue field. At its Frobenius, the eigenvalues on each component Tate module have complex absolute value $q^{1/2}$, by the Weil theorem \cite[II, Theorem 1.1(b)]{MilneAVA567}.

Here is why this verifies the precise associated-weight condition of the cited definition. If a $\Q$-cyclotomic character $\psi$ occurs after passing to an open arithmetic subgroup and satisfies $\psi^b=\chi_\ell^a$, make that further finite extension before spreading. On Frobenius its complex absolute value is both $q^{a/b}$ and $q^{1/2}$. Its associated weight $2a/b$ is therefore $1$. Thus all associated weights of the displayed component quotient lie in $\{1\}\subset(0,2)$, for every characteristic cover. This is exactly $\ell$-graphical fullness. Here $\Q$-cyclotomic is the source's rational-power condition; its distinct term quasi-cyclotomic means agreement with $\chi_\ell$ itself on an open subgroup. The action of the normal tensor changes the gluing, but not these component Tate modules. Corollary 2.7(ii) now applies to the given compatible isomorphism and proves graphicity.
\end{proof}

\begin{lemma}[Known with explicit comparison: normal tensors and twists]\label{A6-normal-kummer}
Fix a stable marked curve $Y/K$ as above. After a finite extension splitting its nodes and branches, choose bases for their normal lines. For two propagators $N=(c_e)$ and $N'=(c'_e)$, geometric transport identifies their admissible groups, and the difference of their arithmetic outer actions is the cocycle
\[
 \bigl(\kappa(c'_e/c_e)\bigr)_e
 \ \in Z^1\!\left(G_K,\prod_e\widehat{\Z}(1)\right).
\]
Changing that transport changes the cocycle by a coboundary. The group of graphic automorphisms that act identically on the underlying graph and on every component anabelioid is precisely the displayed product of profinite Dehn twist groups. This statement includes loops and self-nodes.
\end{lemma}
\begin{proof}
The twist-kernel assertion is exactly \cite[Definition 4.4, Proposition 4.5, Theorem 4.8(iv),(v), and Proposition 4.10(ii), summarized as Theorem B]{HoshiMochizuki2012A567}; it is a theorem for arbitrary semi-graphs of PSC-type, with no tree hypothesis. Definition 4.4 identifies the subgroup trivial on the graph and all component anabelioids; Proposition 4.5 identifies the same condition inside the outer automorphism group. Theorem 4.8(iv) is the product decomposition, and (v) makes the conjugation action the node permutation action times the intrinsic cyclotomic character. Its intrinsic cyclotome is identified with $\widehat{\Z}(1)$ by the geometric orientation.

The same paper's Lemma 5.4(i),(ii) gives local smoothing coordinates $t_e$ and identifies the geometric logarithmic fundamental group of the completed deformation base with the product of their Kummer inertia groups. Its Theorem 5.7 identifies universal monodromy on these factors with the corresponding profinite Dehn twists, including their cyclotomic identifications. Its complete-trait hypothesis is met after base change to $\overline{K}[[t]]$: the residue field is separably closed of characteristic zero, and $\Sigma$ can be all primes. A stable smoothing exists by the local deformation coordinates of Lemma 5.4(i); send its node coordinates to $c_et$. The arithmetic equivariance needed here follows from functoriality of the monodromy map and the twist decomposition. When there is a node, the synchronization in Theorem 5.7 is unique on that factor, so the synchronization commutes with the $G_K$-action as well. When there are no nodes there is no twist factor to compare.

We spell out the arithmetic comparison, rather than infer it from an unframed log base. Over the normal torus the finite Kummer covers have coordinates $u_e^m=t_e$. The tangential points $N,N'$ give fibres $u_e^m=c_e,c'_e$. A geometric path between them is a compatible choice of roots $r_{e,m}$ with
$r_{e,m}^m=c'_e/c_e$. For $\gamma\in G_K$, the ratio
$\gamma(r_{e,m})/r_{e,m}\in\mu_m$ is the Kummer cocycle at level $m$. Inverse limits give the displayed cocycle. By the universal-monodromy comparison just cited, its action on the curve group is the node twist with that same exponent. Any other compatible roots multiply $r_{e,m}$ by a compatible system of roots of unity and hence alter the cocycle by a coboundary.

The nonsmoothing deformation coordinates contribute no extra inertia in Lemma 5.4. Replacing local coordinates $x,y$ by coordinates with derivatives $a,b$ multiplies both the tensor coefficient and the normal coordinate by $ab$; the ratio $c'_e/c_e$ and this comparison are unchanged. Higher-order units have their compatible roots in the strictly henselian formal neighbourhood and have the same effect as a change of geometric path. Thus this calculation depends on the tensors themselves, not their coordinates or an auxiliary smoothing beyond first order. Finally, if a different graphic comparison has the same restrictions to all component anabelioids, its ratio with this one lies in the twist kernel established above. This proves the full stated ambiguity.
\end{proof}

\begin{theorem}[Proved under the specified logarithmic interpretation: 049(a)]\label{A6-anabelian}
Let $K$ be finitely generated over $\Q$, and let $(Y,N)$ and $(Y',N')$ be geometrically connected stable marked curves with nonzero normal tensors as above. For the framed logarithmic fundamental groups, the natural map
\[
 \operatorname{Isom}_K\bigl((Y,N),(Y',N')\bigr)
 \longrightarrow
 \operatorname{Isom}^{\mathrm{out}}_{G_K}
       \bigl(\Delta_{Y,N},\Delta_{Y',N'}\bigr)
\]
is bijective. Labels, if prescribed, must be preserved on both sides. In particular, the graph, component curves, marked points, and actual normal tensors are all recovered. The use of a framed logarithmic group is a hypothesis of this theorem and is the interpretation needed for the manuscript's normal propagator.
\end{theorem}
\begin{proof}
Let $\alpha$ be an isomorphism on the right. Lemma~\ref{A6-graphic} recovers the graph and the component anabelioids. Pass temporarily to a finite extension $L/K$ over which components, nodes, and branches are defined. This restriction is legitimate because $\alpha$ is compatible with the action of all of $G_K$, hence also with $G_L$.

We justify precisely the arithmetic component groups used next. If $\Delta_v\subset\Delta_{Y,N}$ is a chosen vertex subgroup, \cite[Proposition 1.2(ii)]{Mochizuki2007A567} says it is commensurably terminal, hence its geometric normalizer is itself. Since $G_L$ fixes its conjugacy class, its normalizer in the augmented arithmetic group surjects onto $G_L$, with kernel $\Delta_v$. The induced outer action on $\Delta_v$ is the ordinary arithmetic action on the corresponding punctured normalization component: restriction of admissible covers to that component is Galois-equivariant, and changing the normal tensor changes only the gluing. The group $\Delta_v$ is centre-free. An extension with centre-free kernel is the pullback of $\operatorname{Aut}(\Delta_v)\to\operatorname{Out}(\Delta_v)$ along its outer action; this follows directly by sending an element to its conjugation action and its image in the quotient. Thus the recovered normalizer is exactly the augmented arithmetic group of that smooth component, up to the usual geometric inner conjugacy.

The induced isomorphism over $G_L$ is therefore geometric by Mochizuki's smooth anabelian theorem \cite[Theorem A, equivalently Theorem 16.5]{Mochizuki1999A567}. That theorem explicitly covers both full profinite and geometrically pro-$p$ groups; we use its full profinite clause. Its hypotheses apply: these component curves are hyperbolic by stability, and a finitely generated extension of $\Q$ embeds in a finitely generated extension of $\Q_p$, so $L$ is sub-$p$-adic. Applying the theorem to the isomorphism and its inverse, and then using uniqueness for their composites, gives an isomorphism of each smooth component. It extends uniquely to the compact normalization and carries the punctures as dictated by the recovered inertia groups. The graph identifies which pairs are nodes. Gluing these component isomorphisms therefore gives an isomorphism $f:Y_L\to Y'_L$. Uniqueness on components makes it invariant under $\operatorname{Gal}(L/K)$, so it descends to $K$.

Use $f$ to identify the underlying marked curves. Its induced maps on the component groups agree with the maps recovered from $\alpha$. The two geometric comparisons thus differ by an element of the twist kernel in Lemma~\ref{A6-normal-kummer}. The $G_L$-compatibility of $\alpha$ says exactly that the Kummer cocycle of the node-wise tensor ratios is a coboundary. Lemma~\ref{A6-kummer}, applied to $L$, implies that every ratio is $1$. Consequently $f$ preserves $N$ and descends as an isomorphism of pairs.

It remains to show that the originally given $\alpha$ is precisely the outer isomorphism induced by $f$, rather than merely an isomorphism yielding the same curve. Their ratio is a $G_L$-fixed element of the twist kernel. Its node-wise factors are $\widehat{\Z}(1)$, whose fixed parts vanish by Lemma~\ref{A6-kummer}. Thus the ratio is the identity. Conversely an isomorphism of pairs induces the required compatible outer isomorphism by functoriality. If two such isomorphisms induce the same one, smooth anabelian uniqueness on every component shows that they coincide. This proves bijectivity, including the case with no nodes, when the twist kernel is zero.
\end{proof}

\begin{proposition}[Refuted as a bare criterion; corrected recovery in 049(b)]\label{A6-procyclic}
The condition of being isomorphic to $\widehat{\Z}$, with invariant conjugacy class and cyclotomic arithmetic action, does not put such subgroups in bijection with nodes and marked points. Even with all points rational, infinitely many subgroups satisfy it. The strongest reconstruction statement established here is instead the canonical recovery of the full finite graph and all its maximal node and cusp inertia groups in Lemma~\ref{A6-graphic}. It uses the arithmetic action on finite covers, not merely the proposed condition on an individual subgroup.
\end{proposition}
\begin{proof}
Use the split four-pointed nodal curve of Proposition~\ref{A6-ordinary-counterexample}, now with its framed logarithmic group and tensor coefficient $1$. Its group is $\widehat{F}_3$. A peripheral inertia group $I\simeq\widehat{\Z}(1)$ at one of the four rational markings has invariant conjugacy class and cyclotomic action. For every integer $m\ge1$, its characteristic subgroup $mI$ has the same properties and is itself isomorphic to $\widehat{\Z}(1)$. These subgroups are pairwise nonconjugate: in the abelianization $\widehat{\Z}^3$, the chosen inertia has primitive image $\widehat{\Z} e_1$, and $mI$ has image $m\widehat{\Z} e_1$. These closed subgroups are distinct for distinct positive integers $m$, by their prime-adic valuations. Conjugation acts trivially on abelianization. Thus the proposed set is infinite while the set of nodes and markings has five elements.

In this example, the node loop can be taken as the product of the first two peripheral generators and also has primitive nonzero abelianized image. Hence the same defect already affects a node, independently of the cusp example. Taking maximal inertia groups, or identifying commensurable subgroups, removes this particular defect; we do not infer from that repair alone that every maximal invariant procyclic subgroup is geometric. The positive assertion is the reconstruction proved by Lemma~\ref{A6-graphic}, whose derivation supplies the requisite extra information.
\end{proof}

\begin{proposition}[Known: the precise smooth-curve cyclotomic criterion]
\label{A6-cyclotomic-detection}
Let $U/k$ be a smooth affine geometrically connected hyperbolic curve
over a number field, with $\Pi=\pi_1(U)$ and geometric kernel $\Delta$.
For $1\ne z\in\Delta$, the element $z$ belongs to a conjugate of a cusp
inertia group if and only if there is a closed subgroup $N\subset\Pi$
whose image in $G_k$ is open and such that
\[
 nzn^{-1}=z^{\chi(n)}\qquad(n\in N).
\]
Here $\chi$ is the cyclotomic character pulled back from $G_k$.
\end{proposition}
\begin{proof}
This is \cite[Theorem (2.1.1), pp.~94--97]{Nakamura1994A567}.
For a cusp, its decomposition group gives $N$ by the local Kummer
description. Conversely, Lemma~(2.1.4) there says that a nonperipheral
$z$ survives outside the span of cusp classes in the mod-$\ell$
abelianization of some open subgroup $H\subset\Delta$ containing $z$.
After a finite extension of $k$, realize $H$ by an \'etale cover $V$
and make $N\cap\pi_1(V)$ surject onto the base Galois group.
The image of $z$ in the Tate module of the compactification's Jacobian
would then be cyclotomic. At good reduction its Frobenius eigenvalue
would have absolute value $q$, while the Jacobian eigenvalues have
absolute value $\sqrt q$. Thus that image is zero. The kernel is the
span of cusp classes, contradicting the choice of $H$.
\end{proof}
This detects elements belonging to inertia, including their powers;
it does not make all procyclic subgroups into a finite set. Its
hypotheses concern smooth affine curves with their usual arithmetic
action. For the framed singular-logarithmic action, the applicable
statement above is Lemma~\ref{A6-graphic}.

\begin{proposition}[Refuted literally; clarified geometric reading of NB 048]\label{A6-three-classes}
For every hyperbolic type $(g,n)$ in characteristic zero, there are infinitely many conjugacy classes of injective continuous homomorphisms $\widehat{\Z}(1)\to\Delta_{g,n}$. Thus the NB cannot literally assert that there are three such classes, or two when $n=0$.

For geometric simple loops there is a meaningful three-way division into peripheral, separating essential, and nonseparating loops. This is a division into kinds, not three conjugacy classes. Under the mapping class group with unordered punctures, peripheral loops have one orbit when $n>0$, nonseparating loops one orbit when $g>0$, and separating loops have orbits indexed by unordered admissible pairs
\[
 \{(h,r),(g-h,n-r)\},\qquad
 2h+r-1>0,\quad 2(g-h)+(n-r)-1>0.
\]
For labelled punctures their subsets must also be specified.
\end{proposition}
\begin{proof}
The abelianization of the surface group is $\widehat{\Z}^{\,2g+n-1}$ for $n>0$, and $\widehat{\Z}^{\,2g}$ for $n=0$. Its rank is positive in every hyperbolic case. Choose a geometric generator whose abelianized image is primitive. Its procyclic closure is $\widehat{\Z}$: the surjection from $\widehat{\Z}$ onto this closure followed by the primitive abelian coordinate is the identity. The homomorphisms sending $1$ to its distinct positive integral powers are injective and have different abelianized images of $1$, hence are not conjugate.

For the geometric division, cut the oriented surface along a simple loop. If the complement is connected, it has genus $g-1$ and two new boundary circles; classification of oriented surfaces gives one orbit. If it is disconnected, the genera and puncture subsets of the two sides are precisely the complete invariants under homeomorphisms. A side is neither a disc nor a once-punctured disc exactly when its displayed inequality holds. The excluded once-punctured disc is the peripheral case. Gluing boundary-parametrized homeomorphisms of the cut surfaces proves sufficiency of these invariants. This is the change-of-coordinates argument of \cite[\S1.3]{FarbMargalit2012A567}. We do not silently assign this interpretation to the ambiguous NB, nor does it assert that every profinite procyclic subgroup comes from a simple loop.
\end{proof}

\subsection{What the cited combinatorial papers do and do not supply}
The mathematical content suggested by L6 is supported in the following precise senses. In \cite[Corollary 3.11]{Mochizuki2006A567}, for $j=1,2$ let $K_j$ be a finite extension of $\Q_{p_j}$ and let $X_j^{\log}/K_j$ be a smooth log curve. An isomorphism of their \emph{geometric tempered} fundamental groups determines, functorially, an isomorphism of the semi-graphs of anabelioids with compact structure of the geometric special fibres of their stable models; it also forces $p_1=p_2$. This follows there from Theorem 3.7 and Corollary 3.9. The author's later corrections \cite[items (2), (3), (6)]{Mochizuki2020CommentsA567} clarify the proof and require the prime set to contain both residue characteristics in the pro-$\Sigma$ version. That topology must not be silently replaced by a plain profinite group.

The profinite graphical-fullness theorem used here is \cite[Corollary 2.7(ii)]{Mochizuki2007A567}; its hypotheses were verified in Lemma~\ref{A6-graphic}. Hoshi--Mochizuki \cite[Theorem 4.1]{HoshiMochizuki2011A567} gives vertex and node preservation for compatible nodally nondegenerate actions when one vertex group is already matched. Their Corollary 4.2 gives graphicity when there are cusps and the isomorphism is already cuspidal. These are conditional implications, not the bare criterion of 049(b). Theorem~\ref{A6-anabelian} additionally uses the smooth theorem and the 2012 twist and monodromy results; merely listing the 2006 and 2011 papers would not derive it.

The historical links in L6 can also be made more precise from the
now-opened originals. Nakamura \cite[Theorem (5.10), p.~215]{Nakamura1990A567}
proves that, for finite sets $\Lambda,\Lambda'\subset\mathbf P^1(K)$
of cardinality at least three and $K$ finitely generated over $\Q$,
a $G_K$-compatible arithmetic fundamental-group isomorphism determines
a projective transformation carrying $\Lambda$ to $\Lambda'$ with
the same induced puncture bijection. Theorem (6.1) on that page descends
the resulting isomorphism for arbitrary smooth hyperbolic genus-zero
curves. These statements concern reconstruction of the curve; they
do not assert that every given group isomorphism is induced by that
curve isomorphism. Definition (3.3), p.~209, and Theorem (3.4), p.~210,
give an inertia criterion over number fields with an essential
geometric self-normalizing condition. The later, opened 1994 result
supplies the precise element criterion stated above.

Ihara--Nakamura \cite[\S2.3.6, equation $(5_\mu)$, p.~131;
Theorems 1 and $1'$, pp.~132--133]{IharaNakamura1997A567}
construct and algebraize a formal deformation with equations
$TT'=q_{i(\mu)}$. Its inputs include a tangential choice at every
branch and a numbering of the nodes; their $\mathbf P^1_{01\infty}$
diagrams have smooth irreducible components and thus exclude self-nodes.
On the diagonal $q_i=q$, Theorem $2'$ in \S3.6, p.~146, compares
the geometric tangential groupoids of the smooth generic fibre and
the three-punctured line: a Galois element trivial on the latter
pro-$\mathcal C$ groupoid has a compatible lift trivial on the former,
for an almost full class $\mathcal C$ of finite groups.
Theorems 3A--3B and Corollary 4.2.2, p.~148, give the constant-field
inclusion and the pro-$\ell$ equality for $n\ge1$.
The same page reports Takao's removal of that restriction;
that report is distinguished from their proved corollary.
These are deformation and Galois-action comparison results.
The full-faithfulness assertion of Theorem~\ref{A6-anabelian}, including
self-nodes, uses the stronger sources and derivations already specified.
The original articles were recovered from the official GDZ
article links advertised in its volume manifests; their mathematical
pages, rather than only their metadata or the 1995 announcement,
were inspected.

\subsection{049(c): a precise definition, specialization, and a reconstruction criterion}
\begin{proposition}[Proved: the unadorned fundamental group cannot select a compactification]\label{A6-blowup}
Let $U=(\mathbf P^1_\Q\setminus\{0,1,\infty\})^2$. The same smooth variety, and hence the same augmented arithmetic fundamental group, admits two nonisomorphic smooth proper compactifications with strict normal-crossings boundary, containing respectively six and seven irreducible boundary components.
\end{proposition}
\begin{proof}
Take $X=\mathbf P^1_\Q\times\mathbf P^1_\Q$ and the six horizontal and vertical boundary lines. Blow up the rational boundary crossing $(0,0)$, giving $X'$. The blow-up is an isomorphism over $U$. Near the centre the boundary is $xy=0$; the two blow-up charts have coordinates $(x,v)$ with $y=xv$ and $(u,y)$ with $x=uy$. The reduced inverse image of the boundary is respectively $xv=0$ and $uy=0$, hence is again a strict normal-crossings divisor. Its components are the six strict transforms and the exceptional line, seven in total. An isomorphism of pairs preserves the number of irreducible boundary components, so these compactifications are not isomorphic. Their common open variety is literally $U$, and its arithmetic fundamental group and outer $G_\Q$-action are therefore identical.
\end{proof}

The counterexample does not refute the request to \emph{define} an anabelian compactification. The sheet does not prescribe uniqueness, minimality, or a reconstruction axiom. Here is a definition with a proved specialization theorem, followed by an additional condition under which reconstruction is also proved.

\paragraph{A tangential logarithmic definition.}
Fix a field $k$ of characteristic zero and a class $\mathcal A$ of smooth geometrically connected $k$-varieties for which the arithmetic fundamental-group functor is fully faithful on \emph{isomorphisms}. A \emph{tangential log compactification} of $U\in\mathcal A$ consists of a smooth proper pair $(X,D)$, with $D$ a reduced strict normal-crossings divisor and $U=X\setminus D$, together with its boundary-stratum and normal-torus diagrams. Precisely, for every nonempty set $J$ of boundary branches take
\[
 S_J=\left(\bigcap_{j\in J}D_j\right)
       \setminus\bigcup_{i\notin J}D_i,\qquad
 T_J=\prod_{j\in J,S_J}
       \left(N_{D_j/X}|_{S_J}\right)^\times .
\]
One can equivalently use the associated fs divisorial log scheme and its strata. Branch labels, when chosen, and their Galois descent are part of the data. A point of $T_J$ supplies nonzero normal directions; it determines the tangential fibre functor on finite covers of $U$. Locally choose boundary equations $x_j$ and adjoin their compatible roots. A change $x_j\mapsto u_jx_j$ with $u_j|_{S_J}=1$ has compatible roots with residue $1$, so the resulting functors are canonically identified. This is the first-order independence proved in \cite[Propositions 1--2 and Remark 1]{HoshiMatsumotoNakayama2026A567}. The groupoid of these fibre functors, including their isomorphism torsors, retains the normal data that would be discarded by retaining only the bare compact scheme.

The construction globalizes over the whole normal torus without choosing a retraction of a neighbourhood onto the stratum. Pull the log structure of $X$ back along $T_J\to S_J\to X$. On $T_J$ let $v_j$ be its tautological nonzero normal vector. In a local boundary chart put $a_j=dx_j(v_j)$, an invertible function. The elements $x_j/a_j$ of the pulled-back log structure agree on overlaps: if $x'_j=u_jx_j$, then $a'_j=(u_j|_{S_J})a_j$, and that same restricted unit is the transition in the pulled-back log structure. They therefore give a global chart with characteristic monoid $\mathbb N^J$. Adjoin all compatible roots of this chart, exactly as in the source's pro-Kummer construction. Pulling covers back to that pro-log scheme and using its equivalence with the ordinary finite-cover category of $T_J$ gives a functor
\[
 \operatorname{F\acute Et}(U)\longrightarrow
 \operatorname{F\acute Et}(T_J).
\]
Its fibre at a geometric point is the stated tangential fibre functor. The transition equalities satisfy the cocycle condition, so this also descends when geometric boundary branches are permuted by Galois. Composition along nested strata has the canonical homotopy of the source's Proposition 3. Thus this definition specifies a global diagram of finite-cover functors, not merely unrelated pointwise local charts.

This definition permits different models of the same $U$, as Proposition~\ref{A6-blowup} requires. It does not declare that every stratum is itself anabelian, or that a chosen vector in a normal line is reconstructed from the bare group. Those would be further requirements.

\begin{proposition}[Known with derivation: existence in arbitrary dimension]\label{A6-SNC-existence}
Every smooth separated finite-type variety over a characteristic-zero field has a smooth proper compactification with reduced strict normal-crossings boundary. Therefore every member of $\mathcal A$ has a tangential log compactification in the stated sense.
\end{proposition}
\begin{proof}
Apply \cite[Theorem 1.1.15]{Temkin2018A567} with $S=\operatorname{Spec}k$ and with the author's $Y_0=Y=U$, so the initial boundary is empty. A field is excellent; smoothness makes $U$ regular. The theorem gives a regular proper $k$-scheme containing $U$ with strict normal-crossings complement. Over the perfect field $k$, regular finite-type schemes are smooth. Taking the reduced complement gives the asserted pair; its strata, normal tori and tangential functors are then the constructions above.
\end{proof}

The following specialization has no dimension bound.

\begin{theorem}[Known with derivation: logarithmic tangential specialization]\label{A6-log-specialization}
Let $k$ have characteristic zero, $R=k[[t]]$, $F=k((t))$, and let $S^{\log}$ be $\operatorname{Spec}R$ with the divisorial log structure of $t=0$. Let $\mathcal X^{\log}\to S^{\log}$ be a proper morphism of noetherian fs log schemes with $\mathcal X^{\log}$ log regular. Suppose its interior $U$ is contained in the generic fibre and is geometrically connected over $F$. Put
\[
 X_0^{\log}=\mathcal X^{\log}\times_{S^{\log}}s^{\log},
\]
where $s^{\log}$ is the strict closed log point. Restriction gives equivalences
\[
 \operatorname{F\acute Et}(U)
 \simeq\operatorname{K\acute Et}(\mathcal X^{\log})
 \simeq\operatorname{K\acute Et}(X_0^{\log}).
\]
Consequently there is an outer isomorphism
\[
 \pi_1^{\mathrm{et}}(U)\simeq
 \pi_1^{\mathrm{k\acute et}}(X_0^{\log})
 \quad\hbox{over}\quad
 G_F\simeq\pi_1^{\mathrm{k\acute et}}(s^{\log}).
\]
The frame $t$ gives a section
$\sigma_t:G_k\to\pi_1^{\mathrm{k\acute et}}(s^{\log})
\simeq\widehat{\Z}(1)\rtimes G_k$. Pullback along that section defines a tangential specialization extension
\[
 1\longrightarrow\pi_1^{\mathrm{et}}(U_{\overline{F}})
 \longrightarrow
 \pi_1^{\mathrm{k\acute et}}(X_0^{\log})
       \times_{\pi_1^{\mathrm{k\acute et}}(s^{\log})}G_k
 \longrightarrow G_k\longrightarrow1.
\]
In particular this applies to proper log-smooth saturated semistable pairs, with horizontal boundary included. It retains the logarithmic special fibre, rather than replacing it by the ordinary singular fibre.
\end{theorem}
\begin{proof}
The second equivalence is \cite[Corollary 3.9, pp.\ 31--32 of the July 2006 preprint]{HoshiLog2009A567}: properness over a complete local ring and log regularity are its hypotheses. The first is its Proposition A.10, pp.\ 57--58, identifying finite Kummer-étale covers of a log-regular scheme with finite étale covers of its interior tame at the generic points of the boundary. Every such ramification here has residue characteristic zero and is tame. These equivalences commute with pullback from $S^{\log}$. Apply the same two results to $S^{\log}$ itself to identify its group and that of $s^{\log}$ with $G_F$. This proves the compatibility over the base, not just an abstract isomorphism of total groups.

The chart $1\mapsto0$ of the standard log point identifies its geometric Kummer covers with adjoining roots of the chart generator. Their automorphisms form $\widehat{\Z}(1)$; the action on roots of unity is the cyclotomic action. Keeping that generator fixed gives the asserted splitting. The usual arithmetic fundamental-group sequence of the geometrically connected $F$-variety $U$ is exact. Pulling that sequence back along $\sigma_t$ proves the displayed extension. No cohomological comparison theorem is being substituted for the required equivalence of finite-cover categories.

For completeness, the dependence on the frame is explicit. If $t'=at$ and compatible roots of $a\in k^\times$ are chosen, the section fixing the roots of $t'$ is the section fixing the roots of $t$ multiplied by the inverse Kummer cocycle of $a$. Changing $t$ by a factor in $1+tR$ introduces no further ambiguity, since each such unit has a unique $m$th root with constant coefficient $1$ for every $m$. Thus the section depends on the nonzero first-order normal direction. This is the higher-dimensional version of the framing issue in Theorem~\ref{A6-anabelian}; the proof of specialization itself has no dimension bound.
\end{proof}

There is also a geometric form of this assertion. For a proper, separated, exact, relatively smooth map of log analytic spaces with fine base and relatively coherent source, \cite[Theorem 5.1, p.\ 2226]{NakayamaOgus2010A567} makes the Kato--Nakayama realization a topological fibre bundle. A path in the realized base therefore gives an outer isomorphism between the fibre fundamental groups. In a semistable family over a disc, a boundary angle specifies a limiting fibre with the same topology as a nearby smooth fibre. Keeping that angle is essential. This supplies a direct geometric specialization statement; it is not a statement that the ordinary singular fibre has the same group.

\begin{theorem}[Proved: reconstruction of a log-canonically polarized compactification]\label{A6-canonical-compactification}
Let $(X,D)$ and $(X',D')$ be smooth proper pairs over a characteristic-zero field $k$, with reduced strict normal-crossings boundaries. Suppose
$L=\omega_X(D)$ and $L'=\omega_{X'}(D')$ are ample. Every isomorphism
$U=X\setminus D\simeq X'\setminus D'=U'$ extends uniquely to an isomorphism of pairs.

Consequently, on any isomorphism-anabelian class $\mathcal A$, the map
\[
 \operatorname{Isom}_k((X,D),(X',D'))
 \longrightarrow
 \operatorname{Isom}_{G_k}
       \bigl(\pi_1^{\mathrm{et}}(U),\pi_1^{\mathrm{et}}(U')\bigr)
       /\operatorname{Inn}\bigl(\pi_1^{\mathrm{et}}(U'_{\overline{k}})\bigr)
\]
is bijective for these compactifications. In particular their strata, normal-torus diagrams and the associated tangential fibre functors are recovered functorially.
\end{theorem}
\begin{proof}
We first prove independence of the graded log-pluricanonical ring
\[
 R_{\log}(U)=\bigoplus_{m\ge0}
      H^0\bigl(X,\omega_X(D)^{\otimes m}\bigr)
\]
from the chosen smooth proper strict normal-crossings compactification. The birational map between two such models is an isomorphism on $U$. Weak factorization \cite[Theorem 0.3.1, especially (6)]{AKMW2002A567} factors it over $k$ into blow-ups and blow-downs along smooth centres outside $U$ and having normal crossings with the boundary.

It suffices to treat one nontrivial blow-up $f:X_1\to X$ along a connected smooth centre $Z$ of codimension $c\ge2$, contained in exactly $r$ boundary components at its generic point. Let $E$ be its exceptional divisor, and let $D_1$ be the reduced inverse image of $D$. Smoothness and the normal-crossings condition give
\[
 K_{X_1}=f^*K_X+(c-1)E,\qquad
 f^*D=D_{\rm strict}+rE,\qquad
 D_1=D_{\rm strict}+E.
\]
Since $Z$ is contained in the intersection of those $r$ transverse divisors, $r\le c$. Hence
\[
 K_{X_1}+D_1=f^*(K_X+D)+(c-r)E .
\]
For every integer $q\ge0$, $f_*\mathcal O_{X_1}(qE)=\mathcal O_X$: a section is a rational function regular away from the codimension-at-least-two centre, and therefore extends across it because $X$ is normal; conversely every regular function pulls back to a section. The projection formula now gives equality of the spaces of log $m$-canonical forms, under their common realization as rational $m$-canonical forms on $U$. These equalities respect multiplication. Disconnected centres are handled component by component, and blowing up a Cartier divisor is an isomorphism. Iterating proves the ring's independence and functoriality under isomorphisms of $U$.

For an ample line bundle on a proper scheme, the scheme is the Proj of its full section ring; one may use a sufficiently high very ample Veronese to see this directly. Thus
\[
 X\simeq\operatorname{Proj}R_{\log}(U),\qquad
 X'\simeq\operatorname{Proj}R_{\log}(U').
\]
An isomorphism of interiors induces an isomorphism of these graded rings, and hence of their Proj schemes. It extends the given interior map because the induced function-field map is the original one: ratios of equal-degree sections recover the function field for an ample bundle. The complement of $U$ recovers $D$, so the extension preserves the pairs. Separatedness and density prove uniqueness. Finally apply the assumed isomorphism-anabelian theorem to recover the unique map of interiors, then the extension just proved. Normal bundles, strata and their finite-cover fibre functors are functorial constructions from that pair.
\end{proof}

The ample condition gives a concrete choice in Proposition~\ref{A6-blowup}: on $\mathbf P^1\times\mathbf P^1$ the log canonical bundle is $\mathcal O(1,1)$, whereas on the blow-up it is its pullback and has degree zero on the exceptional curve. The latter model is therefore excluded. In dimension one the condition says exactly $2g-2+n>0$, so it recovers the usual hyperbolic completion.

The isomorphism-anabelian hypothesis has substantial concrete realizations. Hoshi's published \cite[Corollary 3.18]{Hoshi2014A567} gives the required bijection for hyperbolic polycurves over sub-$p$-adic fields when either dimension is at most four. Such a polycurve is a tower of hyperbolic curve fibrations; ordinary polycurves are special cases of the orbi-polycurves in that corollary. Carlson's 2026 \emph{preprint} \cite[Theorem A = Theorem 3.11]{Carlson2026A567} asserts the same Isom theorem for all hyperbolic polycurves over a sub-$p$-adic field, without a dimension bound. Its proof applies its Corollary 3.16 and Hoshi's polynondegenerate Hom theorem to both a group isomorphism and its inverse, then uses faithfulness to make the two resulting morphisms inverse. The canonical-compactification theorem above is independent of that recent preprint: it states explicitly the anabelian input it needs.

Thus 049(c) admits a precise answer at the level of a chosen tangential log compactification and a proved specialization theorem, and an additional proved reconstruction criterion when the log canonical bundle is ample. No assertion is made that every smooth anabelian variety admits such an \emph{ample} smooth compactification, that every possible compactification is recoverable, or that a boundary stratum is anabelian without further hypotheses. Those stronger universal existence and reconstruction requirements are not specified on the sheet and are not established by the theorems above. They must remain separate from the definition and specialization that have now been supplied.

\subsection{A compactification system independent of the chosen model}

The definition on 049(c) can retain every smooth normal-crossings model
at once. This gives a choice-independent object even when no individual
model is distinguished by the open variety.

\begin{theorem}[Proved: the canonical system of smooth compactifications]
\label{A6-model-system}
Let $k$ be a field of characteristic zero and let $U/k$ be smooth,
geometrically connected, separated and of finite type. Define
$\mathcal C_k(U)$ as follows. An object is a smooth proper integral
$k$-scheme $X$, with a specified dense open immersion $j_X:U\hookrightarrow X$
whose reduced complement $D_X$ is a strict normal-crossings divisor.
A morphism from $X$ to $Y$ is a $k$-morphism $f:X\to Y$ satisfying
$f\circ j_X=j_Y$. Then:
\begin{enumerate}[label=(\roman*)]
\item $\mathcal C_k(U)$ is essentially small, nonempty and cofiltered.
There is at most one morphism between any two objects. Every morphism
is proper and satisfies $f^{-1}(j_Y(U))=j_X(U)$.
\item Its inverse system, with its marked common interior,
\[
 \mathfrak C_k(U)=\bigl((X,D_X)\bigr)_{X\in\mathcal C_k(U)},
 \qquad U\longrightarrow\mathfrak C_k(U),
\]
is independent, up to canonical pro-isomorphism, of the choice of a
small skeleton. For any fixed model $X_0$, restricting to all models
dominating $X_0$ gives the same pro-object. The same holds after fixing
any finite collection of initial models.
\item Every $k$-isomorphism $\varphi:U\to V$ extends uniquely to an
isomorphism of these pro-systems with their marked interiors. Thus,
on any class $\mathcal A$ for which the augmented arithmetic fundamental
group is fully faithful on isomorphisms, that group reconstructs
$\mathfrak C_k(U)$ functorially. The reconstructed object retains the
schemes, boundaries and transition maps, not just the abstract indexing
partially ordered set.
\end{enumerate}
Here an isomorphism of marked pro-systems includes the isomorphism of
their interiors and commutes with the displayed maps from those
interiors. The theorem makes no selection of one compactification.
\end{theorem}
\begin{proof}
Since $U$ is smooth and geometrically connected, it is geometrically
integral. Nonemptiness is Theorem~\ref{A6-SNC-existence}. Essential
smallness follows because finite-type $k$-schemes, finite affine gluing
data and morphisms between them form a set up to isomorphism; take any
small skeleton when forming the inverse system.

Two morphisms $X\to Y$ agreeing on $U$ are equal: the equalizer is closed
because $Y$ is separated, and equality on the dense open determines a
morphism from the reduced scheme $X$. Every such morphism is proper.
Indeed its graph is closed in $X\times_kY$, whose projection to $Y$ is
proper. Put $W=f^{-1}(j_Y(U))$. The given copy of $U$ in $W$ is a section
of the separated morphism $W\to U$, hence is closed in $W$; it is also
open. The nonempty open $W$ of the integral scheme $X$ is irreducible,
so this open-and-closed copy of $U$ is all of $W$. This proves the asserted
preimage equality. In particular the divisorial log structures of the
pairs are preserved by all transition maps.

For two models $X_1,X_2$, let $\Gamma$ be the reduced closure of the
diagonal copy of $U$ in $X_1\times_kX_2$. It is integral and proper over
$k$. The intersection with $j_{X_1}(U)\times_kX_2$ is exactly the graph
of $j_{X_2}$, because that graph is closed over $U$ by separatedness.
Consequently $U$ is open in $\Gamma$. Let $Z=\Gamma\setminus U$ with
its reduced closed subscheme structure. Every singular point of
$\Gamma$ lies in $Z$. Apply Temkin
\cite[Theorem~1.1.11, p.~4]{Temkin2018A567} to $(\Gamma,Z)$.
Its hypotheses hold: $\Gamma$ is noetherian and quasi-excellent, being
of finite type over the characteristic-zero field $k$. The theorem
gives a modification $X_3\to\Gamma$ supported on $Z$, with $X_3$
regular and with total transform of $Z$ strictly monomial. Its reduced
support is therefore a strict normal-crossings divisor. The modification
is an isomorphism over $U$; $X_3$ is proper over $k$, and regularity over
the perfect field $k$ gives smoothness. The two projections make $X_3$
an object dominating $X_1$ and $X_2$ in $\mathcal C_k(U)$.

This proves the common-domination condition. Parallel arrows are already
equal, so the equalizing condition for a cofiltered category holds as
well. For a fixed $X_0$, every model $X$ has a common refinement with
$X_0$, and any two such refinements have a further common refinement.
Thus the models over $X_0$ are cofinal in the precise inverse-system
sense: the category of their arrows to any given $X$ is nonempty and
connected. Iterating gives the assertion for finitely many fixed models.
The defining formula
\[
 \operatorname{Hom}_{\mathrm{Pro}}
       ((X_i),(Y_j))
   =\varprojlim_j\varinjlim_i\operatorname{Hom}(X_i,Y_j)
\]
then proves that these restrictions and changes of skeleton give
canonical pro-isomorphisms.

Let $\varphi:U\to V$ be an isomorphism. For each model $Y$ of $V$,
regard $Y$ as a model of $U$ using the marking $j_Y\circ\varphi$.
The identity map of that scheme gives a representative of a pro-map
$\mathfrak C_k(U)\to\mathfrak C_k(V)$. These representatives are
compatible with every transition map. Equivalently, to use any prescribed
initial source model, resolve the closure of the graph of $\varphi$ in
its product with $Y$, by the common-domination construction above.
Two representatives extending $\varphi$ agree on a common refinement:
they agree on its dense open $U$ and the target is separated. Hence the
pro-map is unique. Applying the construction to $\varphi^{-1}$ and using
this uniqueness makes both composites the identity. It also proves
compatibility with composition.

For $U,V\in\mathcal A$, a compatible outer isomorphism of augmented
arithmetic groups first yields its unique isomorphism $\varphi:U\to V$
by the defining Isom-anabelian hypothesis. The construction just proved
extends that particular isomorphism uniquely. Conversely an isomorphism
of the marked pro-systems restricts to its isomorphism of interiors,
which induces the original group isomorphism. This is a deduction from
the stated anabelian input; it imposes no anabelian hypothesis on the
individual boundary strata. The strata, their normal bundles and the
normal-torus families of tangential cover functors are reconstructed
for each model because they are intrinsic constructions from that
scheme and its boundary. This recovers the family of possible frames,
not a distinguished frame or point of a normal torus.
\end{proof}

The supported modification used here is provided by Temkin's theorem.
For a description by successive local operations, the birational map
between any two smooth proper models also has the weak factorization
of Abramovich--Karu--Matsuki--W{\l}odarczyk
\cite[Theorem~0.3.1, especially~(6), pp.~2--3]{AKMW2002A567}:
the base field may be any characteristic-zero field, the varieties are
complete and nonsingular, and the centers lie outside $U$ and have
normal crossings with the relevant boundary transforms. This is a
zigzag of blow-ups and blow-downs. The common-domination proof above
requires no assertion that both maps from its common model are
composites of blow-ups along smooth centers.

\begin{theorem}[Proved: compatibility of tangential specialization across models]
\label{A6-model-specialization}
Let $k$ have characteristic zero, put $R=k[[t]]$ and $F=k((t))$, and fix
a smooth geometrically connected separated finite-type $F$-scheme $U$.
Consider proper integral regular $R$-schemes $\mathcal X$ containing
$U$ as a dense open, with reduced complement $D_{\mathcal X}$ a strict
normal-crossings divisor. Give them their divisorial log structures,
and give $S=\operatorname{Spec}R$ the log structure of $t=0$.
Their morphisms are $R$-morphisms restricting to $\operatorname{id}_U$.
This category is nonempty and cofiltered. Every such model is flat
over $R$, and it has a canonical morphism of log schemes to $S^{\log}$.

For every model, write $\mathcal X_0^{\log}$ for its strict closed
log fibre. The equivalences
\[
 \operatorname{F\acute Et}(U)
 \simeq\operatorname{K\acute Et}(\mathcal X^{\log})
 \simeq\operatorname{K\acute Et}(\mathcal X_0^{\log})
\]
are compatible with every transition map. In particular, if
$f:\mathcal X'\to\mathcal X$ is one, then
\[
 f_0^*:\operatorname{K\acute Et}(\mathcal X_0^{\log})
       \longrightarrow
       \operatorname{K\acute Et}(\mathcal X_0'{}^{\log})
\]
is an equivalence. The comparisons between any two models are independent
of common refinement, up to the unique natural isomorphism compatible
with restriction to $U$. They preserve the maps to $G_F$ and, for the
same chosen frame $t$, the tangential specialization extensions over
$G_k$ of Theorem~\ref{A6-log-specialization}.
\end{theorem}
\begin{proof}
The ring $R$ is excellent and noetherian. The map $U\to S$ is separated
and of finite type, since $F=R[1/t]$ is of finite type over $R$.
Apply Temkin \cite[Theorem~1.1.15, p.~5]{Temkin2018A567} with
$Y_0=Y=U$ and with base $S$. It produces a regular proper model with
the required strict normal-crossings complement. Taking the component
containing $U$ removes any extraneous components; in the construction
in \cite[\S3.5.3, p.~24]{Temkin2018A567}, $U$ is already dense in its
compactification. This proves nonemptiness. All the schemes in the
stated category are integral and dominate $S$. Their coordinate rings
are torsion-free $R$-modules and hence are flat over the discrete
valuation ring $R$. The function $t$ has divisor supported on the
boundary. Locally it is a unit times a product of powers of boundary
parameters, so it defines the indicated morphism of divisorial log
schemes over $S^{\log}$.

The common-domination proof of Theorem~\ref{A6-model-system} works over
$R$: take the closure of the graph of $U$ in
$\mathcal X_1\times_R\mathcal X_2$ and apply
\cite[Theorem~1.1.11]{Temkin2018A567} to its closed complement of $U$.
This graph closure is a noetherian quasi-excellent characteristic-zero
scheme, because it is of finite type over the excellent ring $R$.
Its resolution is regular, proper and integral over $R$, is unchanged
on $U$, and has strict normal-crossings boundary. It is again flat by
the torsion-free argument. Uniqueness of arrows and the equality
$f^{-1}(U)=U$ follow from the same density and section argument as
before. Thus all arrows induce canonical log morphisms over $S^{\log}$.

A regular scheme with its strict normal-crossings divisorial log
structure is log regular. The cover equivalences are therefore exactly
the application of Hoshi's
\cite[Proposition~A.10 and Corollary~3.9]{HoshiLog2009A567}
proved in Theorem~\ref{A6-log-specialization}: every residue
characteristic is zero, so boundary ramification of a finite cover of
$U$ is tame, and properness over the complete local base supplies the
closed-fibre equivalence. All these hypotheses hold for every model
in the present category; semistability or a saturated log-smooth
structure is not being silently assumed.

For an arrow $f$, restriction and pullback give the commuting diagram
\[
 \begin{array}{ccc}
 \operatorname{K\acute Et}(\mathcal X^{\log})&
   \xrightarrow{\ f^*\ }&
 \operatorname{K\acute Et}(\mathcal X'{}^{\log})\\
 \downarrow&&\downarrow\\
 \operatorname{F\acute Et}(U)&\xrightarrow{\ \operatorname{id}\ }&
 \operatorname{F\acute Et}(U).
 \end{array}
\]
The vertical functors are equivalences, so $f^*$ is one. The analogous
square of restrictions to the strict closed log fibres then shows that
$f_0^*$ is an equivalence. More precisely, if $j_{\mathcal X}$ and
$i_{\mathcal X}$ denote the interior and closed-fibre inclusions, equip
each special-fibre cover category with the comparison
\[
 E_{\mathcal X}=j_{\mathcal X}^*\circ(i_{\mathcal X}^*)^{-1}
       :\operatorname{K\acute Et}(\mathcal X_0^{\log})
        \longrightarrow\operatorname{F\acute Et}(U),
\]
including the natural isomorphisms of the chosen quasi-inverse.
Every comparison between models is equipped with its isomorphism to
the identity under these $E_{\mathcal X}$. Full faithfulness supplies
a unique natural isomorphism between any two comparisons that respects
these specified identifications. Bare uniqueness without that
normalization is not asserted. Consequently the comparisons
are independent of refinement, and their isomorphisms obey the cocycle
condition. This is the stated categorical meaning of canonicity.

All functors in the two squares commute with pullback from $S^{\log}$.
Thus their induced outer fundamental-group isomorphisms commute with
the augmentations to $G_F$. Keeping the same section
$\sigma_t:G_k\to G_F$ and pulling back these augmented groups gives
the same tangential extension over $G_k$. Changing the frame has the
Kummer dependence already computed in Theorem~\ref{A6-log-specialization};
changing the model introduces no additional ambiguity. The construction
is for models of this fixed generic open over this fixed trait. It
does not assert that an unrelated compactification over $k$ is a model
of that family.
\end{proof}

\begin{proposition}[Refuted: reconstruction of a chosen special fibre]
\label{A6-special-fibre-counterexample}
Over $R=\Q[[t]]$, there are two proper integral regular models of the
same smooth geometrically connected generic open, both with strict
normal-crossings reduced boundary, whose special fibres are not
isomorphic. Nevertheless their strict special log fibres have equivalent
finite Kummer-\'etale cover categories over $G_{\Q((t))}$ and the same
tangential arithmetic extension over $G_\Q$ for the fixed frame $t$.
Thus even these augmented and framed cover data do not reconstruct a
chosen special-fibre scheme. This does not contradict reconstruction of
the complete marked model system.
\end{proposition}
\begin{proof}
Put $F=\Q((t))$, take
\[
 \mathcal X=\mathbf P^1_R\times_R\mathbf P^1_R,
 \qquad
 U=(\mathbf P^1_F\setminus\{0,1,\infty\})^2,
\]
and let $D$ be the union of the six horizontal boundary divisors and
the vertical fibre $\mathcal X_0$. This is a strict normal-crossings
divisor on the proper regular integral $R$-scheme $\mathcal X$.
Blow up the closed rational point $P=(x,y,t)=(0,0,0)$ and write
\[
 f:\mathcal X' =\operatorname{Bl}_P\mathcal X\longrightarrow\mathcal X,
 \qquad D'=(f^{-1}D)_{\mathrm{red}}.
\]
The blow-up is projective, is an isomorphism away from $P$, and fixes
$U$. Near $P$ the three boundary parameters are $x,y,t$. Its charts are
\[
 \begin{array}{c|c|c|c}
 \text{chart}&\text{substitution}&D'&t\text{ on the chart}\\\hline
 x&y=xv,\ t=xu&xvu=0&xu\\
 y&x=ya,\ t=yb&yab=0&yb\\
 t&x=tc,\ y=td&tcd=0&t .
 \end{array}
\]
The blow-up of the regular parameter ideal $(x,y,t)$ is regular:
the completed local chart equations $t-xu$, respectively $t-yb$,
eliminate $t$ and leave regular power-series rings; the $t$-chart is
regular directly. Away from the exceptional divisor the map is
an isomorphism. The vanishing factors in the table form part of a
regular parameter system, proving the strict normal-crossings
condition. The other horizontal divisors miss $P$ and are unchanged
near the exceptional locus. The total space $\mathcal X'$ is integral
and dominates $\operatorname{Spec}R$, so it is flat over this discrete
valuation ring. Both models therefore satisfy every hypothesis of
Theorem~\ref{A6-model-specialization}.

The old special fibre is the irreducible surface
$\mathcal X_0=\mathbf P^1_\Q\times\mathbf P^1_\Q$.
Since $t$ vanishes to order one at $P$, the new special fibre is reduced
and is the union
\[
 \mathcal X'_0=
 \operatorname{Bl}_{(0,0)}
       (\mathbf P^1_\Q\times\mathbf P^1_\Q)
 \ \cup_{\mathbf P^1_\Q}\ \mathbf P^2_\Q .
\]
The second component is the exceptional divisor
$\mathbf P(T_P\mathcal X)$; its intersection with the first is
$\mathbf P(T_P\mathcal X_0)$. The two factors $xu$, respectively $yb$,
in the table verify multiplicity one for both components. Thus the
special fibres have respectively one and two geometric irreducible
components and cannot be isomorphic, even after extending $\Q$.

Give both total spaces their boundary log structures and both special
fibres their strict induced log structures. The cited theorem identifies
their Kummer-\'etale cover categories by $f_0^*$, with the identifications
normalized by their common restriction to $\operatorname{F\acute Et}(U)$.
These identifications commute with pullback from the standard log trait.
Consequently their fundamental groups are isomorphic over the same
$G_F$, up to the usual basepoint conjugacy. Pulling back along the same
root-fixing section $\sigma_t:G_\Q\to G_F$ gives identical augmented
extensions with geometric kernel $\pi_1(U_{\overline F})$ over $G_\Q$.
In particular the associated tangential outer actions agree, although
the special-fibre schemes differ. The algebraic model and blow-up also
descend to $\Q[t]_{(t)}$; the generic open over $\Q(t)$ is a
dimension-two hyperbolic polycurve. No unproved anabelianity hypothesis
is needed for this counterexample or for its cover comparison.
\end{proof}


\subsection{Presentations on 67/050--051}
We write $\operatorname{Mod}(S_{g,n})$ for orientation-preserving mapping classes preserving the set of punctures, $\operatorname{PMod}$ for those fixing them individually, and $\operatorname{Mod}^{\pm}$ for the extended group. A group fixing one distinguished puncture may still permute the others. The original image 051 has a clear $+!$ in its first formula, but the superscript in the second formula can be read as $+'$; we keep the two group interpretations separate.

\begin{theorem}[Proved: the two matrix presentations on 050]\label{A6-matrix-presentations}
The presentations, with the involution relations made explicit, are
\[
 \operatorname{SL}_2(\Z)
 =\langle \sigma,\rho\mid \sigma^2=\rho^3,\ \rho^6=1\rangle,
\]
\[
 \operatorname{GL}_2(\Z)
 =\langle r_1,r_2,r_3\mid
 r_1^2=r_2^2=r_3^2=1,\
 (r_1r_2)^2=(r_2r_3)^3,\
 (r_1r_2)^4=1\rangle.
\]
We may take
\[
 \sigma=\begin{pmatrix}0&-1\\1&0\end{pmatrix},\quad
 \rho=\begin{pmatrix}0&-1\\1&1\end{pmatrix},\quad
 r_1=\begin{pmatrix}-1&0\\0&1\end{pmatrix},\
 r_2=\begin{pmatrix}0&1\\1&0\end{pmatrix},\
 r_3=\begin{pmatrix}1&1\\0&-1\end{pmatrix}.
\]
In particular these identify the ordinary and extended once-punctured torus groups.
\end{theorem}
\begin{proof}
The projective modular group is $C_2*C_3$, with the images of $\sigma,\rho$ as generators; see \cite[Appendix C, Theorem C.1]{ConradSL2A567}. In the first abstract presentation, $z=\sigma^2=\rho^3$ is central and $z^2=1$. The quotient by $z$ is exactly $C_2*C_3$. The indicated matrix homomorphism is surjective: $\sigma^{-1}\rho=\left(\begin{smallmatrix}1&1\\0&1\end{smallmatrix}\right)$ and the Euclidean algorithm reduces any unimodular integer column using this shear and $\sigma$. Its kernel is contained in $\langle z\rangle$, and $z$ maps to $-I\ne I$. Thus it is injective.

In the second presentation put $\sigma=r_1r_2$, $\rho=r_2r_3$, $r=r_2$. It becomes the first presentation together with
\[
 r^2=1,\qquad r\sigma r=\sigma^{-1},\qquad r\rho r=\rho^{-1}.
\]
The given $r_2$ induces exactly this automorphism of $\operatorname{SL}_2(\Z)$. Hence the abstract group is $\operatorname{SL}_2(\Z)\rtimes C_2$. The determinant map on $\operatorname{GL}_2(\Z)$ has this kernel and is split by $r_2$, proving the second presentation. Direct matrix multiplication, including all defining relations, is checked in the attached exact program. The identification of torus mapping classes with their actions on $H_1$ is \cite[Theorem 2.5 and \S2.2.4]{FarbMargalit2012A567}; its punctured version fixes the origin.
\end{proof}

\begin{proposition}[Proved: the four-punctured sphere and both superscript readings]\label{A6-four-punctures}
For the four-punctured sphere,
\[
 \operatorname{PMod}(S_{0,4})\simeq F_2
 \simeq \Gamma(2)/\{\pm I\},
 \qquad
 \operatorname{Mod}(S_{0,4})
 \simeq V_4\rtimes\operatorname{PSL}_2(\Z).
\]
The subgroup fixing one specified puncture is $\operatorname{PSL}_2(\Z)$, and the quotient of the full group by its displayed normal $V_4$ is also $\operatorname{PSL}_2(\Z)$. Thus the formula with $+!$ is false if that symbol consistently means the pure group; it is correct if the second superscript means the distinguished-point subgroup. The exact diagram intended by the main part of 051 is
\[
 \begin{array}{ccccccccc}
 1&\longrightarrow&F_2&\longrightarrow&
 \operatorname{Mod}(S_{0,4})&\longrightarrow&S_4&\longrightarrow&1\\
 &&\Vert&&\downarrow&&\downarrow\\
 1&\longrightarrow&\Gamma(2)/\{\pm I\}&\longrightarrow&
 \operatorname{PSL}_2(\Z)&\longrightarrow&S_3&\longrightarrow&1 .
 \end{array}
\]
\end{proposition}
\begin{proof}
The torus quotient by $z\mapsto-z$ is the sphere with the images of its four two-torsion points marked. A mapping class lifts to the torus in two ways, giving its matrix modulo $\pm I$. The kernel consists of translations by two-torsion. Conversely linear maps give a section. This is the precise lifting argument of \cite[Proposition 2.7]{FarbMargalit2012A567}, and yields the semidirect product. On the four two-torsion points the action is the affine group
$\mathbf F_2^2\rtimes\operatorname{GL}_2(\mathbf F_2)\simeq S_4$.
Its linear subgroup fixes zero and is the full stabilizer of that point; the same holds upstairs, proving the distinguished-point assertion. Its pure kernel is therefore the kernel of reduction modulo $2$ in $\operatorname{PSL}_2(\Z)$, namely $\Gamma(2)/\{\pm I\}$.

Independently, forgetting the fourth labelled puncture gives the Birman sequence with base the triply punctured sphere. Its pure mapping class group is trivial and its fundamental group is $F_2$, so the pure four-point group is $F_2$. This proves the identification and the left vertical equality in the diagram. Finally $F_2$ is torsion-free, whereas $\operatorname{PSL}_2(\Z)=C_2*C_3$ contains elements of orders two and three; hence the two conflicting literal readings cannot define the same group.
\end{proof}

\begin{proposition}[Proved: a consistent reading of the partly erased lift on 051]\label{A6-binary-S3}
There is an exact sequence
\[
 1\longrightarrow F_2\longrightarrow\operatorname{SL}_2(\Z)
 \longrightarrow\widetilde{S}_3\longrightarrow1,
\]
where $\widetilde{S}_3$ is the binary dihedral group of order $12$. It can be obtained as a quotient of $\operatorname{SL}_2(\Z/4\Z)$, as suggested by the partially erased formula on the sheet. The quotient by $\{\pm I\}$ gives the lower row of Proposition~\ref{A6-four-punctures}. This is not the ordinary reduction-modulo-$2$ sequence, whose kernel contains $-I$.
\end{proposition}
\begin{proof}
In $\operatorname{SL}_2(\Z/4\Z)$ put
\[
 W=\left\{
 I,\begin{pmatrix}3&2\\0&3\end{pmatrix},
 \begin{pmatrix}3&0\\2&3\end{pmatrix},
 \begin{pmatrix}1&2\\2&1\end{pmatrix}\right\}.
\]
These are $I+2A$, where $A$ runs through the two-dimensional subspace
$\{0,I+E_{12},I+E_{21},E_{12}+E_{21}\}$ of the trace-zero matrices over $\mathbf F_2$. Conjugation by the reductions of $\sigma$ and the elementary shear permutes the three nonzero elements, so $W$ is normal. It is a Klein four group and does not contain $-I$. The full kernel of reduction from level $4$ to level $2$ has eight elements, so it is $W\times\{\pm I\}$. The preimage $K$ of $W$ in $\operatorname{SL}_2(\Z)$ is consequently a complement to $\{\pm I\}$ in $\Gamma(2)$ and maps isomorphically to $\Gamma(2)/\{\pm I\}=F_2$. The quotient has order $48/4=12$.

In that quotient the images of $\rho,\sigma$ satisfy
$\rho^6=1$, $\sigma^2=\rho^3$, and
$\sigma\rho\sigma^{-1}=\rho^{-1}$. The first two matrices retain orders six and four because $-I\notin W$. These relations describe the binary dihedral group of order twelve, and they generate the quotient. Enumeration of all $48$ unimodular matrices modulo $4$, the normal subgroup $W$, and these quotient relations is included in the exact program. The image supports this possible reading of $\widetilde{S}_3$; it does not resolve every overwritten symbol on the sheet.
\end{proof}

\begin{proposition}[Proved: the integral projective-line quotient on 051]\label{A6-projective-line}
With
\[
 B_\infty=\left\{\ \pm\begin{pmatrix}1&n\\0&1\end{pmatrix}
                    :n\in\Z\right\},
\]
there are bijections of $\operatorname{SL}_2(\Z)$-sets
\[
 \operatorname{SL}_2(\Z)/B_\infty
 \simeq\mathbf P^1(\Z)\simeq\mathbf P^1(\Q).
\]
Reduction gives the displayed surjection
$\mathbf P^1(\Z)\to\mathbf P^1(\mathbf F_2)$.
These are point sets; the quotient written on the sheet is not an identification with the entire scheme $\mathbf P^1_{\Z}$.
\end{proposition}
\begin{proof}
Every rational line has a primitive integral generator $(a,c)$, unique up to sign. Conversely primitive pairs modulo sign are exactly the integral points of the projective line: over $\Z$, their common ideal is the unit ideal and every invertible module is free. Bezout supplies $b,d$ with $ad-bc=1$, so $(a,c)$ is the first column of an element of $\operatorname{SL}_2(\Z)$. The stabilizer of the line generated by $(1,0)$ is precisely $B_\infty$, proving the two bijections. A primitive pair is never $(0,0)$ modulo $2$; changing its sign has no effect on the resulting projective point. The pairs $(1,0)$, $(0,1)$ and $(1,1)$ lift the three target points, proving surjectivity. Finally $\mathbf F_2^\times=1$, so $\operatorname{GL}_2(\mathbf F_2)=\operatorname{SL}_2(\mathbf F_2)=\operatorname{PGL}_2(\mathbf F_2)\simeq S_3$, explaining the interchangeable finite-group labels in that part of the sheet.
\end{proof}

The peripheral relations written nearby have the usual oriented interpretation. A four-punctured sphere has presentation
$\langle\ell_0,\ell_1,\ell_2,\ell_3\mid
\ell_0\ell_1\ell_2\ell_3=1\rangle$.
Cutting it into two pairs of pants gives
$\ell_0\ell_1c=1$ and $c'\ell_2\ell_3=1$, with $c'=c^{-1}$ when both boundary orientations are induced from their respective pants. Eliminating $c,c'$ gives the four-puncture relation. An equality $c=c'$ instead uses opposite orientation conventions on the two copies; it must not be combined with two identical induced-boundary conventions.

\begin{proposition}[Known with the missing hypotheses: the Birman sequence]\label{A6-birman}
Let $S=S_{g,n-1}$ be an oriented surface with $2g-2+(n-1)>0$, and let $p$ be an additional marked point. There are exact sequences
\[
 1\longrightarrow\pi_1(S,p)\longrightarrow
 \operatorname{Mod}(S,p)\longrightarrow\operatorname{Mod}(S)
 \longrightarrow1,
\]
and their pure and extended versions. Here the middle group fixes $p$; forgetting one point is not a map from a group that is allowed to interchange that point with all the others. The same form holds with boundary fixed pointwise. The positivity condition is essential.
\end{proposition}
\begin{proof}
Apply \cite[Theorem 4.6 and the remarks immediately following it]{FarbMargalit2012A567}. Its hypothesis is exactly $\chi(S)<0$, and its evaluation-fibration proof identifies the kernel with point pushing. More explicitly, the relevant portion of the homotopy sequence is
\[
 \pi_1(\operatorname{Homeo}_0^+(S))\longrightarrow
 \pi_1(S,p)\longrightarrow
 \pi_0(\operatorname{Homeo}^+(S,p))
 \longrightarrow \pi_0(\operatorname{Homeo}^+(S))
 \longrightarrow1.
\]
For negative Euler characteristic the first contribution vanishes, giving the assertion. Restrictions to pure groups and allowance of orientation reversal give the stated versions.
For the closed torus, translations make the evaluation map surjective on $\pi_1(S)$, and $\operatorname{Mod}(S_{1,1})\to\operatorname{Mod}(S_{1,0})$ is the isomorphism $\operatorname{SL}_2(\Z)\to\operatorname{SL}_2(\Z)$. Thus the written left injection would be false there. The twice-punctured sphere gives another failure when a third puncture is added: its $\pi_1$ is $\Z$ but both pure mapping class groups are trivial. These are precisely why the hyperbolicity hypothesis belongs to the corrected sequence.
\end{proof}

\begin{proposition}[Known with derivation: pants, Lego, and involutions]\label{A6-presentation-types}
The three presentation mechanisms anticipated on 050 have rigorous realizations, with different conventions. The pants and Lego mechanisms give complete mapping-class \emph{groupoid} presentations by local moves and their relations. Closed extended mapping class groups of every genus $g\ge1$ admit finite presentations on orientation-reversing involutions; generation by three such involutions is known. These assertions do not by themselves identify an arithmetic automorphism group of the full Teichm\"uller tower.
\end{proposition}
\begin{proof}
Bakalov--Kirillov work with compact oriented surfaces with boundary and a marked point on each boundary circle. Their Theorem 4.24 \cite{BakalovKirillov2000A567} states that the complex of markings, with the $Z,F,B,S$ moves and the relations of their \S\S4.7 and 4.22, is connected and simply connected. Their Theorem 5.1 gives the equivalent presentation restricted to maximal markings, that is, pants decompositions (with the specified elementary pieces and their markings). Join any two markings by an edge path; connectivity supplies one, and simple connectivity says that any two are related by the listed two-cell relations. Thus these moves and relations completely present the marking-change groupoid.

The distinction emphasized in their Remark 4.3 matters: an edge is not in general a homeomorphism, and the simply connected marking-change groupoid by itself has trivial vertex automorphism groups. To recover the mapping class group one retains its action on the marking complex and takes the resulting action groupoid, with stabilizers retained. Explicitly, fix a marking $M$; the pairs $(g,[p])$, where $g$ is a mapping class and $p$ is a path from $M$ to $gM$, multiply by $(g,[p])(h,[q])=(gh,[p\cdot gq])$. Connectivity and simple connectivity give exactly one path class for each $g$, so this group is the mapping class group. Passing to representatives of marking orbits expresses it by the local moves, piece automorphisms, and the indicated two-cell relations, while stabilizers must not be discarded as they would be in an ordinary topological quotient. This is the precise action-groupoid sense of the pants/Lego presentation here. Capping a boundary circle by a once-marked disc quotients the corresponding boundary twist, so puncture and boundary conventions must be compared by that quotient, not identified without it.

For involutions, Stukow's \cite[Theorem 3.1]{Stukow2004A567} proves that the extended mapping class group of a closed orientable genus-$g$ surface, $g\ge1$, is generated by three orientation-reversing involutions. Finite presentability of the ordinary mapping class group is proved in \cite[\S5.3]{FarbMargalit2012A567}; adjoining one orientation-reversing involution and its conjugation relations gives a finite presentation of the extended group. A finitely presented group has a finite presentation on any prescribed finite generating set: express the old generators as words in the new ones, introduce both sets with those defining equations, and eliminate the old generators by Tietze transformations. Applied to Stukow's three generators, this gives a finite presentation whose generators are the required involutions, with their square relations included. This proves the existence assertion; Stukow's generation theorem alone is not being called a printed list of complete relations. For genus one the explicit complete presentation is Theorem~\ref{A6-matrix-presentations}.

The boundary and punctured forms of the pants/Lego theorem are covered by its stated surface conventions and the capping operation. The three-involution theorem just used has closed-surface scope; we do not extend that numerical bound to all boundary or labelled-puncture conventions. Finally the geometric groupoid theorem describes moves and their homotopies. An arithmetic realization additionally needs compatibility with Galois actions; that information is not a consequence of simple connectivity alone.
\end{proof}

